% \subsection{Probability of Sample Linear Independence}
% Let $v_1 \in \Z_q^t$ be a fixed nonzero vector, and let
% $v_2 \in \Z_q^t$ be sampled uniformly at random.
% The probability that $v_2$ is linearly independent of $v_1$ is
% \[
% \mathbb{P}\bigl(v_2 \notin \langle v_1 \rangle\bigr)
% = \frac{q^t - q}{q^t}
% = 1 - q^{1-t}.
% \]

% Conditioned on $v_1$ and $v_2$ being linearly independent, let
% $v_3 \in \Z_q^t$ be sampled uniformly at random.
% The probability that $v_3$ is linearly independent of $\{v_1,v_2\}$ is
% \[
% \mathbb{P}\bigl(v_3 \notin \langle v_1,v_2 \rangle \,\big|\, v_1,v_2 \text{ independent}\bigr)
% = \frac{q^t - q^2}{q^t}
% = 1 - q^{2-t}.
% \]

% More generally, for $j \geq 2$, conditioned on $v_1,\dots,v_{j-1}$ being linearly independent, a uniformly random vector
% $v_j \in \Z_q^t$ satisfies
% \[
% \mathbb{P}\bigl(v_j \notin \langle v_1,\dots,v_{j-1} \rangle \,\big|\, v_1,\dots,v_{j-1} \text{ independent}\bigr)
% = 1 - q^{j-1-t}.
% \]

% Hence, the probability that $k$ uniformly random vectors
% $v_1,\dots,v_k \in \Z_q^t$ are linearly independent is
% \begin{equation}\label{independence}
% \mathbb{P}_k
% = \prod_{j=1}^{k} \left(1 - \frac{q^{j-1}}{q^t}\right)
% = \prod_{j=1}^{k} \left(1 - q^{j-1-t}\right).
% \end{equation}

% %%%%%%%%%%%%%%%%%%%%%%%%%%%%%%%%

% \noindent
% \textbf{Example.}
% In our setting,
% \[
% t = \binom{n+d}{d}, \qquad
% k = \binom{n+d-1}{d}.
% \]
% Using the CRYSTALS-Kyber parameter set $n=256$, $d=7$, and $q=3329$ (Kyber768, see~\cite{NISTPQC:CRYSTALS-KYBER22}), we obtain
% \begin{itemize}
%     \item $k = \binom{262}{7} \approx 1.55 \times 10^{13}$,
%     \item $t = \binom{263}{7} \approx 1.59 \times 10^{13}$,
%     \item $t - k \approx 4.24 \times 10^{11}$.
% \end{itemize}

% The largest term $q^{j-1-t}$ in~\eqref{independence} occurs for $j = k$, yielding
% \[
% q^{k-1-t} = 3329^{-(t-k+1)} \approx 3329^{-4.24 \times 10^{11}}
% \approx 10^{-10^{12}},
% \]
% which is negligibly small.
% For such values, we may use the approximation $\log(1-x) = -x + \mathcal{O}(x^2)$, giving
% \[
% \log \mathbb{P}_k
% = \sum_{j=1}^{k} \log(1 - q^{j-1-t})
% = -\sum_{j=1}^{k} q^{j-1-t} + \mathcal{O}\!\left(\sum_{j=1}^{k} q^{2(j-1-t)}\right).
% \]
% Exponentiating and retaining the leading term (using $e^x = 1 + x + \mathcal{O}(x^2)$) yields
% \[
% \mathbb{P}_k
% = 1 - \sum_{j=1}^{k} q^{j-1-t} + o(1).
% \]

% Using the bound $\sum_{j=1}^{k} q^{j-1-t} \le k\, q^{k-1-t}$, we obtain
% \[
% \mathbb{P}_k
% \ge 1 - k\, q^{k-1-t}
% \approx 1 - 10^{13} \cdot 10^{-10^{12}}
% > 1 - 10^{10^{2}}\cdot 10^{-10^{12}}
% \approx 1-10^{-10^{10}}.
% \]
% Thus, the probability that the sampled vectors are linearly independent is overwhelmingly close to $1$.


%%%%%%%%%%%%%%%%%%%%%%%%%%%%%%%%
%%%%%%%%%%%%%%%%%%%%%%%%%%%%%%%%


% For small $q^{j-1-m}$ we can use Taylor expansions
% \[
% \mathbb{P}_k = e^{\log(\mathbb{P}_k)} = e^{\sum_{j=1}^k(\log(1-q^{j-1-t}))} =  e^{\sum_{j=1}^k(-q^{j-1-t})} \sim 1 - \sum_{j=1}^k(q^{j-1-t}),
% \]
% where we used: $\log(1-x) = -x +\mathcal{O}(x^2)$ and $e^x=1+x+\mathcal{O}(x^2)$.

% This correspnds to a probability of:
% \begin{align*}
%     \mathbb{P}_k &= 1 - \sum_{j=1}^k(-q^{j-1-t}) \geq \\
%     &\geq 1 - k(q^{k-1-t}) \sim \\
%     &\sim 1 - (10^{13})\cdot \big(10^{-10^{12}}\big) \geq \\
%     &\geq 1 - (10^{10^{2}})\cdot \big(10^{-10^{12}}\big)\sim 1-10^{-10^{10}}.
% \end{align*}


%%%%%%%%%%%%%%%%%%%%%%%%%%%%%%%%
%%%%%%%%%%%%%%%%%%%%%%%%%%%%%%%%


% \subsection{Probability of Diagonal Pivots}
% The probability that $k := \binom{n+d-1}{d}$ vectors
% $v_1,\dots,v_k \in \Z_q^k$ are linearly independent is
% \begin{equation}\label{pivots}
%     \mathbb{P}_k(\{v_1,\dots,v_k\} \text{ independent})
%     = \prod_{j=1}^{k} \left(1 - \frac{q^{j-1}}{q^k}\right)
%     = \prod_{j=1}^{k} \left(1 - q^{j-1-k}\right).
% \end{equation}
% To estimate $\mathbb{P}_t$, we use
% \(\mathbb{P}_t = e^{\log(\mathbb{P}_t)}\) and the inequality
% \(\log(1-x) \ge - \frac{x}{1-x}\) for \(0 < x < 1\):
% \begin{align*}3
% \log \mathbb{P}_t 
% &= \sum_{j=1}^{t} \log\left(1 - q^{j-1-t}\right)\geq \sum_{j=1}^{t} \left(1 - \frac{1}{1 - q^{j-1-t}}\right)\\
% &= - \sum_{j=1}^{t} \frac{q^{j-1-t}}{1 - q^{j-1-t}} \geq - \frac{1}{1 - q^{-1}} \sum_{j=1}^{t} q^{j-1-t}\quad (1 - q^{-1} = \min_{1\le j\le k}{\{1-q^{j-1-t}\}})\\
% &= - \frac{1}{1 - q^{-1}} \frac{\frac{1}{q} - \frac{1}{q^{t+1}}}{1 - \frac{1}{q}} = - \frac{1}{1 - q^{-1}} \cdot \frac{1 - q^{-t}}{q-1} \sim - \frac{1}{(1 - q^{-1})(q-1)}.
% \end{align*}

% \noindent
% \textbf{Example.}  
% Using CRYSTALS-Kyber~\cite{NISTPQC:CRYSTALS-KYBER22}  parameters $q=3329$, we obtain
% \begin{itemize}
%     \item $1 - q^{-1} \approx 0.99969961$;
%     \item $\dfrac{1 - q^{-t}}{q-1} \approx \dfrac{1}{q-1} \approx 3.0048 \cdot 10^{-4}$;
%     \item $\log \mathbb{P}_t \approx - \frac{3.0048 \cdot 10^{-4}}{0.9997} \approx -3.0057 \cdot 10^{-4}$.
% \end{itemize}
% Hence
% \[
% \mathbb{P}_t \ge \exp(-3.0057 \cdot 10^{-4}) \approx 0.99969943,
% \]
% so the probability that the pivots lie on the diagonal is overwhelmingly close to 1.


% KEEP ESTIMATES FOR CLARITY?
% \begin{remark}
%     We will use the following relations to give a bound for $\log(\mathbb{P}_k)$:
%     \begin{itemize}
%         \item $\log x \ge 1-\frac{1}{x}$ for $0<x<1$;
%         \item \item $q^{j-1-t} \le q^{-1}$;
%         \item $\sum_{k=0}^{r}x^k = \frac{1-x^{r+1}}{1-x}$;
%         \item $e^x > 1 + x$ for $x>0$.
%     \end{itemize}
% \end{remark}
% \begin{itemize}
%     \item $1-q^{-1} \approx 0.99969961$;
%     \item $\frac{1 - q^{-t}}{q-1} \approx \frac{1}{q-1} \approx 3.0048 \cdot 10^{-4}$;
%     \item $\log \mathbb{P}_t \geq - \frac{3.0048 \cdot 10^{-4}}{0.9997} \approx -3.0057 \cdot 10^{-4}$.
% \end{itemize}

%%%%%%%%%%%%%%%%%%%%%%%%%%%%%%%%
%%%%%%%%%%%%%%%%%%%%%%%%%%%%%%%%


% \subsection{Duality from~\cite{EC:NMSU25} (Keep it?)}

% \begin{definition}[Learning With Bounded Errors]
%     Let $\mathbb{F}$ be a field.
%     Let $k,n,d\in\mathbb{N}$, an choose subsets $S_1,\dots,S_n\subset \mathbb{F}$, each of size $d$.

%     For a generator matrix $\vars{G}\in\mathbb{F}^{n\times k}$, the LWBE problem asks to extract $\vars{x}\in\mathbb{F}^{k}$ from $(\vars{G},\vars{b}=\vars{Gx+e})$, where $\vars{e}\in S_1\times \cdots\times S_n$.
% \end{definition}

% It is easy to see that LWBE problem is equivalent to its dual version:

% \begin{definition}[Bounded Syndrome Decoding]
%     For a parity-check matrix $\vars{H}\in\mathbb{F}^{(n-k)\times n}$, the BSD problem asks to extract $\vars{e}$ from $(\vars{H},\vars{s}=\vars{He})$, where $\vars{e}\in S_1\times \cdots\times S_n$.
% \end{definition}
% \newpage






\subsection{Probability of Diagonal Pivots}
\luc{This part has to aligned to the provided definition parameters (k should not be used)}

Let \(n < B\) and consider \(n\) vectors \(v_1,\dots,v_n \in \Z_q^B\). 
The probability that these vectors are linearly independent is well-known and can be expressed as
\begin{equation}\label{pivots}
\mathbb{P}_n(\{v_1,\dots,v_n\} \text{ are independent}) 
= \prod_{j=1}^{n} \left(1 - q^{j-1-B}\right).
\end{equation}
For large \(q\), this probability is extremely close to 1.  

\noindent
\textbf{Example.} 
Consider the parameters of CRYSTALS-Kyber768~\cite{NISTPQC:CRYSTALS-KYBER22} with \(n = 256\), \(d = 7\), and \(q = 3329\). 
The initial set of LWE samples can be seen as \(m\) vectors in \(\Z_q^B\), with $m = \binom{n+d-1}{d} = \binom{262}{7}$, and \(B = \binom{n+d}{d} = \binom{263}{7} \).
Using standard estimates for the product in the above formula, we obtain
\begin{equation*}
\mathbb{P}_m \ge \exp\left(-\frac{1}{(1-q^{-1})(q-1)}\right) \approx 0.9997.
\end{equation*}
Hence, the probability that the pivots lie on the diagonal is overwhelmingly close to 1. 
The same argument applies for each $1 \le i \le d$, with $i$ replacing $d$ in the definitions of $m$ and $B$.

\noindent
\textbf{Example FIXED.} 
Consider the parameters of CRYSTALS-Kyber768~\cite{NISTPQC:CRYSTALS-KYBER22} with \(n = 3\), \(N = 256\), \(d = 5\), and \(q = 3329\). 
The initial set of MLWE samples can be seen as \(m\) vectors in \(\Z_q^B\), with $m = \binom{Nn+d-1}{d}/N = \binom{772}{7}/256$, and \(B = \binom{Nn+d}{d} = \binom{773}{7} \).
Using standard estimates for the product in the above formula, we obtain
\begin{equation*}
\mathbb{P}_m \ge \exp\left(-\frac{1}{(1-q^{-1})(q-1)}\right) \approx 0.9997.
\end{equation*}
Hence, the probability that the pivots lie on the diagonal is overwhelmingly close to 1. 
The same argument applies for each $1 \le i \le d$, with $i$ replacing $d$ in the definitions of $m$ and $B$.